\section{无遗憾是可以做到的}
\label{sec:17-Constrained-Guarantees/04-Online-Learning-and-Regret/02}

手上有八个信号源，每天各自给出一个``拦截／放行''的建议，你必须每天选一个照做，第二天才知道谁对谁错。八个源的质量都不稳定：某一段时间 3 号最准，换个季度变成 6 号，没有哪一个长期占优。

最省事的做法是每天跟着``到目前为止累计错得最少''的那一个走。这个策略有个具体的破法：对手只要盯住你此刻跟随的那个源，专门在它身上做手脚，就能让你天天跟错。八个源里只要有一个每天都对，你却可能连续 \(T\) 天全错——追随领先者的遗憾可以线性增长。

问题出在``选一个''这个动作太硬。上一节的承诺看起来谦逊，可它要求对所有对手序列都成立，而任何确定性的、把全部权重压在一个候选上的规则，都能被针对性地击穿。修补的方向由此确定：不要选一个，要按权重混合，而且混合要留有余地。

\subsection{指数权重把``跟随领先者''软化}

Hedge 算法把每个候选 \(i\) 的累计损失 \(L_{i,t}=\sum_{s\le t}\ell_s(i)\) 记下来，第 \(t\) 轮按

\[
p_{i,t}\;\propto\;\exp(-\eta L_{i,t-1})
\]

的概率随机选一个候选，或者在损失是线性的场合直接按这组权重做加权组合。参数 \(\eta>0\) 控制软硬程度：\(\eta\to\infty\) 退化成追随领先者，\(\eta\to0\) 退化成均匀乱选。中间那一档才有保证。

假设每轮损失取值在 \([0,1]\) 内，则

\[
R_T\;\le\;\sqrt{\tfrac{T}{2}\ln N},
\qquad \eta=\sqrt{\tfrac{8\ln N}{T}}.
\]

证明骨架只有三行，用一个势函数把两侧夹住。令 \(W_t=\sum_{i=1}^{N}e^{-\eta L_{i,t}}\)，取 \(\Phi_t=\ln W_t\)，起点是 \(\Phi_0=\ln N\)。一步的比值恰好是一个期望：

\[
\frac{W_t}{W_{t-1}}=\sum_{i}p_{i,t}\,e^{-\eta\ell_t(i)}.
\]

对取值在 \([0,1]\) 的随机变量用 Hoeffding 引理（\hyperref[sec:05-Probability/07-Transforms-and-Bounds/04]{``Chernoff 方法与 Hoeffding 集中界''}里那条对矩母函数的控制），得 \(\ln\E e^{-\eta X}\le-\eta\E X+\eta^2/8\)。逐轮累加：

\[
\begin{aligned}
\Phi_T-\Phi_0
&\le-\eta\sum_{t=1}^{T}\ip{p_t}{\ell_t}+\frac{\eta^2T}{8},\\
\Phi_T&\ge\ln e^{-\eta L_{i^\star,T}}=-\eta L_{i^\star,T}.
\end{aligned}
\]

上式给出 \(\Phi_T\) 的上界，下式给出下界——\(W_T\) 是一堆正数之和，扔掉其余项只留最优候选那一项即可。两边合并、除以 \(\eta\)，就得到

\[
\sum_{t=1}^{T}\ip{p_t}{\ell_t}-L_{i^\star,T}\;\le\;\frac{\ln N}{\eta}+\frac{\eta T}{8},
\]

对 \(\eta\) 求极小便是上面那个界。

\subsection{\texorpdfstring{\(\ln N\) 是这个界最值得看的地方}{ln N 是这个界最值得看的地方}}

候选数只以对数形式进入。\(N\) 从 \(10\) 涨到 \(10^{6}\)，\(\ln N\) 从 \(2.3\) 涨到 \(13.8\)，界只放大 \(\sqrt{13.8/2.3}\approx2.4\) 倍。候选多了十万倍，要多付的不到三倍。

这与统计学习那边的联合界是同一个机制：\hyperref[sec:13-Machine-Learning/10-Learning-Theory/02]{``从联合界到一致收敛''}里，有限假设类的偏差也只以 \(\sqrt{\ln|\mathcal H|/n}\) 进入。指数权重和联合界都在做同一件事——把``有很多候选''这件事的开销压成对数。区别在于那边需要 i.\,i.\,d.\ 假设才能谈``偏差''，这边什么都不需要。

\begin{center}
\begin{tikzpicture}[x=1cm,y=1cm,font=\small,>=stealth]
  \draw[->,black!60] (0,0) -- (7.6,0) node[right,font=\footnotesize] {\(t\)};
  \draw[->,black!60] (0,0) -- (0,4.7) node[above,font=\footnotesize] {累计损失};
  \draw[black!30] (0,0) -- (1,0.62) -- (2,1.24) -- (3,1.86) -- (4,2.48) -- (5,3.10) -- (6,3.72) -- (7,4.34);
  \draw[black!30] (0,0) -- (1,0.35) -- (2,0.80) -- (3,1.40) -- (4,2.10) -- (5,2.85) -- (6,3.60) -- (7,4.20);
  \draw[orange!70!black] (0,0) -- (1,0.30) -- (2,0.60) -- (3,0.90) -- (4,1.20) -- (5,2.10) -- (6,3.00) -- (7,3.90);
  \draw[green!45!black] (0,0) -- (1,0.45) -- (2,0.90) -- (3,1.35) -- (4,1.80) -- (5,2.25) -- (6,2.70) -- (7,3.15);
  \draw[blue!65!black,very thick]
    (0,0) -- (1,0.34) -- (2,0.66) -- (3,0.98) -- (4,1.30) -- (5,2.24) -- (6,2.92) -- (7,3.40);
  \draw[<->,black!70] (7,3.15) -- (7,3.40);
  \node[anchor=west,font=\footnotesize,black!75] at (7.15,3.28) {\(R_T\)};
  \foreach \x in {1,3,5,7} \draw[black!45] (\x,0) -- (\x,-0.1) node[below,font=\scriptsize] {\x};
  \node[anchor=west,align=left,font=\footnotesize,black!70] at (0.35,4.15)
    {切换处有一小段滞后；\\ \(N\) 增大时它只按 \(\sqrt{\ln N}\) 变长};
  \draw[blue!65!black,very thick] (4.55,1.40) -- (5.0,1.40);
  \node[anchor=west,font=\footnotesize] at (5.07,1.40) {Hedge 加权组合};
  \draw[orange!70!black] (4.55,0.92) -- (5.0,0.92);
  \node[anchor=west,font=\footnotesize] at (5.07,0.92) {前期最好的候选};
  \draw[green!45!black] (4.55,0.44) -- (5.0,0.44);
  \node[anchor=west,font=\footnotesize] at (5.07,0.44) {后期最好的候选};
\end{tikzpicture}

\emph{加权组合那条粗线始终贴着当时最低的一条走，只在领先者换人时落后一小段——候选再多，落后的那一段也只按 \(\ln N\) 的平方根变长。}
\end{center}

图里 Hedge 那条线在第 5 轮附近离开橙线、转去贴绿线，中间那一小段就是权重重新分配所需的时间。它不可能为零：要认出橙线不行了，必须先看到它连续吃亏，而这段观察本身就是遗憾的来源。

条件对账。损失有界是硬条件——若某一轮的损失可以任意大，指数权重会被单点炸掉：一个 \(\ell_t(i)=10^{6}\) 直接把候选 \(i\) 的权重打到浮点下溢，此后再也回不来，而它可能只是那一轮运气差。实现里通常先把损失截断或归一化到 \([0,1]\)，这一步不是数值技巧，是恢复定理前提。\(\eta\) 依赖 \(T\) 是第二个条件：最优的 \(\eta\) 里带着 \(T\)，事先不知道总轮数就取不到。补救办法是倍增技巧——把时间轴切成长度 \(1,2,4,8,\dots\) 的段，每段按该段长度重设 \(\eta\) 并重启，界只多一个小常数因子；或者直接用 \(\eta_t=\sqrt{8\ln N/t}\) 的递减序列。第三个条件容易被忘：界里的 \(\sum_t\ip{p_t}{\ell_t}\) 是期望损失，若真的按 \(p_t\) 随机抽一个候选来执行，实际损失只在高概率意义上接近它，需要再加一个鞅集中不等式。

\subsection{在线梯度下降是同一个更新式换了分析目标}

候选是有限个时用指数权重，决策集是连续的凸集时用另一条路。设 \(\mathcal X\subset\R^d\) 凸紧、直径为 \(D\)，每轮的 \(\ell_t\) 凸且次梯度范数不超过 \(G\)。在线梯度下降的更新是

\[
x_{t+1}=\Pi_{\mathcal X}\bigl(x_t-\eta_t g_t\bigr),\qquad g_t\in\partial\ell_t(x_t),
\]

其中 \(\Pi_{\mathcal X}\) 是到凸集的投影。取 \(\eta_t=D/(G\sqrt t)\)，可得 \(R_T=O(DG\sqrt T)\)。

势函数换成``离比较对象有多远''。对任意固定的 \(x\in\mathcal X\)，投影到凸集不会拉远距离（这是投影的非扩张性，见\hyperref[sec:03-Linear-Algebra/03-Orthogonality/02]{``投影与最近点：从一条直线走向一般子空间''}），所以

\[
\norm{x_{t+1}-x}^2\le\norm{x_t-x}^2-2\eta_t\ip{g_t}{x_t-x}+\eta_t^2\norm{g_t}^2 .
\]

整理出内积，再用凸函数的切平面是全局下界这一条（\hyperref[sec:09-Convex-Optimization/03-Convex-Functions/02]{``切平面为何是全局下界''}）把内积换成函数值之差：

\[
\ell_t(x_t)-\ell_t(x)\;\le\;\ip{g_t}{x_t-x}\;\le\;\frac{\norm{x_t-x}^2-\norm{x_{t+1}-x}^2}{2\eta_t}+\frac{\eta_t G^2}{2}.
\]

对 \(t\) 求和，左边正是遗憾。右边第一项在 \(\eta_t\) 不变时望远镜式地相消，只剩 \(D^2/(2\eta)\)；\(\eta_t\) 递减时多出的部分仍被 \(D^2/(2\eta_T)\) 控住。第二项是 \(\frac{G^2}{2}\sum_t\eta_t\)，与 \(\sqrt T\) 同阶。两项一个随 \(\eta\) 减小而变大、一个随 \(\eta\) 减小而变小，\(\eta\sim1/\sqrt t\) 是它们的平衡点。

更新式与\hyperref[sec:09-Convex-Optimization/06-Optimization-Algorithms/01]{``梯度下降真正难在步长''}里那个一模一样，两处分析的目标却完全不同。那里的问题是``离最优值还有多远''，所以需要目标函数在迭代过程中保持不变，需要 \(L\)-光滑来保证一步不会跨过去，步长上界写作 \(2/L\)。这里的目标函数每轮都换，谈``最优值''没有意义，所以比的是累计损失，步长由 \(D\) 与 \(G\) 定而不由曲率定。同一行代码，两套记账方式。

分布假设的处理方式也翻了个个儿。批量分析里，随机梯度要求样本 i.\,i.\,d.\ 才能保证 \(\E g_t=\nabla f(x_t)\)，噪声怎样影响收敛是\hyperref[sec:09-Convex-Optimization/08-Nonconvex-and-Stochastic-Optimization/04]{``SGD 在期望中下降但噪声会留下平台''}的主题。在线分析里根本没有``真梯度''这个概念要去逼近，\(g_t\) 就是第 \(t\) 轮那个损失的次梯度，是什么就是什么。不假设分布，反倒不需要 i.\,i.\,d.

这两套记账之间有一座桥。若 \(\ell_t\) 确实是 i.\,i.\,d.\ 抽样得到的，把在线算法产生的 \(x_1,\dots,x_T\) 取平均，用 Jensen 不等式（\hyperref[sec:05-Probability/07-Transforms-and-Bounds/03]{``Jensen 不等式与随机凸性''}）与遗憾界一起，可以推出这个平均点的超额风险是 \(O(1/\sqrt T)\)。在线到批量的这条转换说明：无遗憾是比 i.\,i.\,d.\ 泛化更强的要求，做到前者就免费得到后者，反过来不成立。

\subsection{\texorpdfstring{\(\sqrt T\) 绕不过去}{根号 T 绕不过去}}

两个算法都给出 \(\sqrt T\) 量级的界。这不是分析不够精细留下的余量——\(\Omega(\sqrt T)\) 是所有算法的共同下界。

论证短得可以完整写出。取 \(N=2\) 个候选，对手不看我们的动作，每轮独立抛一枚公平硬币：正面则候选 1 损失 \(0\)、候选 2 损失 \(1\)，反面反过来。无论算法怎么选，每轮的期望损失都是 \(1/2\)，所以 \(\E\sum_t\ell_t(x_t)=T/2\)，与算法无关。而事后最优候选的累计损失是

\[
\min_{i\in\{1,2\}}L_{i,T}=\frac{T}{2}-\frac{1}{2}\Bigl|\sum_{t=1}^{T}\varepsilon_t\Bigr|,
\qquad \varepsilon_t=\pm1 \text{ 等概率}.
\]

中心极限定理（\hyperref[sec:05-Probability/08-Limit-Theorems/03]{``中心极限定理与 Gaussian 近似''}）给出 \(\E\abs{\sum_t\varepsilon_t}\approx\sqrt{2T/\pi}\)。两者相减，任何算法的期望遗憾都至少是 \(\Theta(\sqrt T)\)。

值得把这个论证的结构看清楚：遗憾之所以必然增长，不是因为算法笨，而是因为\emph{随机波动本身会制造一个事后看起来很聪明的候选}。事后最优是从 \(N\) 个候选里挑出来的，挑选这个动作把噪声的一侧系统性地放大了——这与\hyperref[sec:13-Machine-Learning/10-Learning-Theory/01]{``训练误差为何总是乐观''}里那一百万枚硬币是同一件事，只是那里叫``训练误差偏乐观''，这里叫``事后基准偏强''。既然基准里含着 \(\Theta(\sqrt T)\) 的运气成分，任何算法都不可能追平。

于是 \(O(\sqrt T)\) 与 \(\Omega(\sqrt T)\) 合上了。这是一条真正的不可能结果：不是当前技术做不到，是没有算法能做到。一般 \(N\) 时下界是 \(\Omega(\sqrt{T\ln N})\)，与 Hedge 的上界只差常数。这一部分反复出现的正是这类结论——先说清想要什么，再证明它的上限在哪里，最后确认上限已经被够到。

到这里为止，每轮结束后我们都能看到所有候选各自的损失，哪怕当时没选它。这份信息是免费拿到的，两个算法都用了它：Hedge 要更新全部 \(N\) 个权重，在线梯度下降要算当前点上的次梯度而不只是函数值。真实的决策往往没有这份便利——推荐了这条视频就看不到另一条的点击率，做了这个定价就不知道另一个定价的成交量。下一节把这一条信息通道掐掉，看看要多付多少。
